\section{Native interrupt and relational contract composition}
\label{sec:a14-composition}

\textbf{Question and design.}
This deterministic experiment connects the native interruption mechanism and
the relational contract in one scripted execution. It tests whether an
application adapter can preserve the reviewed-candidate relation through
approval, handoff, rejection and continuation. It uses LangGraph 1.0.8,
langgraph-checkpoint 4.0.0, langchain-core 1.2.13 and SQLAlchemy 2.0.23 with
Python 3.11. The existing \texttt{stores.py}, its source manifest and nineteen
contract source files are copied without edits. Native framework code is also
unmodified. The protocol, runner, independent scorer, scorer controls and
dependencies are frozen as 27 files before two final executions, each with
fresh SQLite databases. The development round remains separate.

\textbf{Adapter responsibilities.}
An author-written graph first pauses over an already registered immutable
candidate proposing \texttt{machine proposal} for the address field. The
scripted reviewer can deny or authorize \texttt{reviewed correction}; only
an affirmative response to the current native interrupt creates the
application grant handle. A second interrupt is an instrumentation barrier
for injecting the controlled handoff, not another approval request.
The graph invokes the unchanged relational gate with a complete two-field
successor. After rejection, a recovery interrupt receives a scripted choice.
Reuse submits the old candidate/grant; fresh authorization visits a new
review interrupt before creating a grant for the replacement. The unchanged
batch field supplies a source-preservation check. X3 assumes the trusted,
content-immutable candidate storage required in main-text Section 3.2. It
tests the resulting identity relations at commit; it does not independently
probe same-ID content mutation or object replacement between check and use.

The graph, payload interpretation, grant-handle registry, interventions and
recording are adapter code. InMemorySaver and pre-commit grant handles are
in process. The database atomically persists the admitted successor and its
authorization/transition relations. This design does not provide or test
restart-safe approval storage, authenticated review or a distributed
transaction between checkpoint and business storage.

\begin{table}[htbp]
\centering\small
\caption{All ten composition conditions, with identical semantic observations in both final executions. Counts are per execution. ``No write'' refers to business versions and transactional relations, not candidate registration.}
\begin{tabularx}{\textwidth}{@{}>{\raggedright\arraybackslash}p{0.29\textwidth}Xr@{}}
\toprule
Condition & Observed behavior & Address commits\\
\midrule
Corrected approval & Current approval commits corrected value with reviewed identity. & 1\\
Denial & No grant handle or successor write. & 0\\
Identity substitution; reuse & First attempt rejects; old reviewed candidate/grant commits. & 1\\
Identity substitution; fresh review & First attempt rejects; old interrupt ID cannot approve replacement; current ID permits it. & 1\\
Wrong executed value & First attempt rejects; original correction commits on reuse. & 1\\
Stale predecessor & Independent batch update persists; stale address attempt rejects. & 0\\
Rollback after compare-and-swap & No partial transactional rows; retry commits once. & 1\\
Rollback before commit & No partial transactional rows; retry commits once. & 1\\
Duplicate native resume & Completed execution has no additional write. & 1\\
Consumed application-grant replay & Direct replay has no additional write. & 1\\
\bottomrule
\end{tabularx}
\end{table}

\textbf{Independent readback.}
Every graph invocation and gate attempt retains a SQLite copy and the native
pending interrupt identifiers, payloads and scripted input. The scorer
imports neither the gate nor LangGraph. It opens fresh read-only SQLite
connections and checks version counts, corrected values, the retained batch
source, and joins among candidate certificates, authorization bindings,
human-decision rows and fact transitions. It also compares the complete
record-version, transition, decision and authorization-binding tables before
and after every rejected or rolled-back attempt. Gate-return status is not
used to determine whether a business effect occurred. Candidate registration
is permitted before rejection and is kept separate from transactional
business effects.

Both executions satisfy all ten conditions. Each contains fourteen gate
attempts: eight admitted corrections and six rejected or rolled-back attempts
with no partial transactional rows. The stale case contributes one separate
independent batch update and no address correction. Repeating a completed
native resume or explicitly replaying a consumed application grant adds no
business version. In the fresh-review case, the new native interrupt ID
differs from the original; submitting the original ID leaves the current
review pending and creates no new grant or write. Submitting the current ID
permits the replacement's corrected transition with a matching authorization.

Four negative controls mutate database copies: an incorrect final value,
a deleted transition, an incorrect authorized value and an extra version.
The independent scorer detects all four. These checks exercise the oracle;
they are not additional task successes or security-attack trials.

\textbf{Interpretation and artifacts.}
The evidence establishes compatibility of native interrupt routing and the
existing relational checks for the specified adapter and scripted paths.
It does not establish autonomous recovery, production performance, concurrent
replay resistance, restart durability, human review validity or protection
against hostile mutation of pending graph state. Two executions test
deterministic repeatability on the same machine, not statistical robustness.
The agent trajectories reported elsewhere remain a separate evidence layer.

\texttt{Native-contract-reproduction.zip} contains \texttt{PROTOCOL.md},
\texttt{FREEZE.json}, \texttt{run.py}, \texttt{score.py},
\texttt{test\_scorer.py}, the unchanged backend, requirements, both final
run directories and all SQLite snapshots. \texttt{REPRODUCE.md} gives
read-only rescoring and fresh-directory rerun commands. No provider access
or model calls are needed. Development records are retained separately in
the local experiment directory and are not pooled with final observations.
